% GENERATED FILE. Edit canonical lessons, then run npm run generate:books.
% canonical-source-hash: fnv1a64:34802e29c8633b55
% canonical-lessons: ML-C172-vevikkuka, ML-C172-kotukkuka, ML-C172-kontuvaruka, ML-C172-vilikkuka, ML-C172-karayuka

\chapter{Cook, Give, Bring, Call, Cry}
\label{ch:ml-cook-give-bring-call-cry}
\hlchaptermodality{172}

\begin{chapteropening}
\textbf{By the end of this chapter:} \emph{I can say to cook, to give, to bring, to call and to cry.}

Complete the last lesson of chapter 172: i can say to cook, to give, to bring, to call and to cry.
\end{chapteropening}

\section[\texorpdfstring{\ml{വേവിക്കുക} (vēvikkuka)}{വേവിക്കുക (vēvikkuka)}]{\ml{വേവിക്കുക} (vēvikkuka) — to cook}

\label{lesson:ML-C172-vevikkuka}



\begin{quote}
Before the new one: say the Malayalam for to bathe, then the Malayalam for to wash.

Say \textquotedblleft{}in front of\textquotedblright{} and \textquotedblleft{}behind.\textquotedblright{} (\emph{Munnil}, \emph{pinnil}.)
\end{quote}

\subsection*{You'll want to know: \ml{വേവിക്കുക}}
\textbf{\ml{വേവിക്കുക}} — \emph{vēvikkuka} — \textquotedblleft{}to cook\textquotedblright{}.

\begin{grammarlens}[title={What you've built}]
One more word for this chapter.
\end{grammarlens}

\subsection*{Your turn}
\begin{itemize}
\raggedright
  \item \emph{Say it:} \emph{vēvikkuka}
  \item \emph{Say it:} \emph{vēvikkuka}, once more
  \item \emph{Say it:} say \emph{vēvikkuka}
  \item \emph{Recall:} say the Malayalam for to bathe, then the Malayalam for to wash, then say \emph{vēvikkuka} again
\end{itemize}

\begin{tcolorbox}[breakable,colback=teal!4,colframe=teal!35!black,title={Before you move on}]
What does \ml{വേവിക്കുക} mean? (\textquotedblleft{}To cook\textquotedblright{}.) Say it once more.
\end{tcolorbox}
\section[\texorpdfstring{\ml{കൊടുക്കുക} (koṭukkuka)}{കൊടുക്കുക (koṭukkuka)}]{\ml{കൊടുക്കുക} (koṭukkuka) — to give}

\label{lesson:ML-C172-kotukkuka}



\begin{quote}
Before the new one: say the Malayalam for to wash, then the Malayalam for to cook.
\end{quote}

\subsection*{You'll want to know: \ml{കൊടുക്കുക}}
\textbf{\ml{കൊടുക്കുക}} — \emph{koṭukkuka} — \textquotedblleft{}to give\textquotedblright{}.

\begin{grammarlens}[title={What you've built}]
One more word for this chapter.
\end{grammarlens}

\subsection*{Your turn}
\begin{itemize}
\raggedright
  \item \emph{Say it:} \emph{koṭukkuka}
  \item \emph{Say it:} \emph{koṭukkuka}, once more
  \item \emph{Say it:} say \emph{koṭukkuka}
  \item \emph{Recall:} say the Malayalam for to wash, then the Malayalam for to cook, then say \emph{koṭukkuka} again
\end{itemize}

\begin{tcolorbox}[breakable,colback=teal!4,colframe=teal!35!black,title={Before you move on}]
What does \ml{കൊടുക്കുക} mean? (\textquotedblleft{}To give\textquotedblright{}.) Say it once more.
\end{tcolorbox}
\section[\texorpdfstring{\ml{കൊണ്ടുവരുക} (koṇṭuvaruka)}{കൊണ്ടുവരുക (koṇṭuvaruka)}]{\ml{കൊണ്ടുവരുക} (koṇṭuvaruka) — to bring}

\label{lesson:ML-C172-kontuvaruka}



\begin{quote}
Before the new one: say the Malayalam for to cook, then the Malayalam for to give.

Say \textquotedblleft{}with\textquotedblright{} two ways. (\emph{Oppaṁ}, \emph{kūṭe}.)
\end{quote}

\subsection*{You'll want to know: \ml{കൊണ്ടുവരുക}}
\textbf{\ml{കൊണ്ടുവരുക}} — \emph{koṇṭuvaruka} — \textquotedblleft{}to bring\textquotedblright{}.

\begin{grammarlens}[title={What you've built}]
One more word for this chapter.
\end{grammarlens}

\subsection*{Your turn}
\begin{itemize}
\raggedright
  \item \emph{Say it:} \emph{koṇṭuvaruka}
  \item \emph{Say it:} \emph{koṇṭuvaruka}, once more
  \item \emph{Say it:} say \emph{koṇṭuvaruka}
  \item \emph{Recall:} say the Malayalam for to cook, then the Malayalam for to give, then say \emph{koṇṭuvaruka} again
\end{itemize}

\begin{tcolorbox}[breakable,colback=teal!4,colframe=teal!35!black,title={Before you move on}]
What does \ml{കൊണ്ടുവരുക} mean? (\textquotedblleft{}To bring\textquotedblright{}.) Say it once more.
\end{tcolorbox}
\section[\texorpdfstring{\ml{വിളിക്കുക} (viḷikkuka)}{വിളിക്കുക (viḷikkuka)}]{\ml{വിളിക്കുക} (viḷikkuka) — to call}

\label{lesson:ML-C172-vilikkuka}



\begin{quote}
Before the new one: say the Malayalam for to give, then the Malayalam for to bring.
\end{quote}

\subsection*{You'll want to know: \ml{വിളിക്കുക}}
\textbf{\ml{വിളിക്കുക}} — \emph{viḷikkuka} — \textquotedblleft{}to call\textquotedblright{}.

\begin{grammarlens}[title={What you've built}]
One more word for this chapter.
\end{grammarlens}

\subsection*{Your turn}
\begin{itemize}
\raggedright
  \item \emph{Say it:} \emph{viḷikkuka}
  \item \emph{Say it:} \emph{viḷikkuka}, once more
  \item \emph{Say it:} say \emph{viḷikkuka}
  \item \emph{Recall:} say the Malayalam for to give, then the Malayalam for to bring, then say \emph{viḷikkuka} again
\end{itemize}

\begin{tcolorbox}[breakable,colback=teal!4,colframe=teal!35!black,title={Before you move on}]
What does \ml{വിളിക്കുക} mean? (\textquotedblleft{}To call\textquotedblright{}.) Say it once more.
\end{tcolorbox}
\section[\texorpdfstring{\ml{കരയുക} (karayuka)}{കരയുക (karayuka)}]{\ml{കരയുക} (karayuka) — to cry}

\label{lesson:ML-C172-karayuka}



\begin{quote}
Before the new one: say the Malayalam for to bring, then the Malayalam for to call.
\end{quote}

\subsection*{You'll want to know: \ml{കരയുക}}
\textbf{\ml{കരയുക}} — \emph{karayuka} — \textquotedblleft{}to cry\textquotedblright{}.

\begin{grammarlens}[title={What you've built}]
One more word for this chapter.
\end{grammarlens}

\subsection*{Your turn}
\begin{itemize}
\raggedright
  \item \emph{Say it:} \emph{karayuka}
  \item \emph{Say it:} \emph{karayuka}, once more
  \item \emph{Say it:} say \emph{karayuka}
  \item \emph{Recall:} say the Malayalam for to bring, then the Malayalam for to call, then say \emph{karayuka} again
\end{itemize}

\begin{tcolorbox}[breakable,colback=teal!4,colframe=teal!35!black,title={Before you move on}]
What does \ml{കരയുക} mean? (\textquotedblleft{}To cry\textquotedblright{}.) Say it once more.
\end{tcolorbox}
